\chapter{\IfLanguageName{dutch}{Inleiding}{Introduction}}
\label{ch:inleiding}

In de hedendaagse samenleving spelen mediabedrijven een cruciale rol in het bieden van informatie. Het is essentieel dat deze informatie op een objectieve wijze wordt gedeeld, maar de realiteit blijkt niet altijd aan onze verwachtingen te voldoen. Het onpartijdigheidsrapport van de \gls{vrm} toont aan dat VTM Nieuws, onderdeel van DPG Media, bepaalde politieke partijen meer spreektijd geeft dan andere \autocite{VanAelst2024}. Tegelijkertijd heeft \gls{hln}, ook onderdeel van DPG Media, volgens dagelijkse rapporten gepubliceerd door het \textcite{CIM2025} al jaren het grootste bereik onder de Vlaamse nieuwswebsites.

Dit onderzoek richt zich op het ontwikkelen van een browserextensie die nieuwsartikelen van \gls{hln} automatisch analyseert en classificeert op basis van journalistieke partijdigheid. De extensie is specifiek gericht op lezers van \gls{hln} die twijfelen aan de neutraliteit van de berichtgeving en die zich bewust willen worden van de mogelijke vooringenomenheid in hun dagelijkse nieuwsconsumptie.Het vergroten van transparantie over mogelijke bias kan bijdragen aan kritisch mediagebruik en het versterken van het vertrouwen in nieuwsmedia.

\section{\IfLanguageName{dutch}{Probleemstelling}{Problem Statement}}
\label{sec:probleemstelling}

Nieuws dat niet op neutrale wijze wordt weergegeven, kan bijdragen aan een afnemend vertrouwen in de traditionele media \autocite{Hanjia2024}. Daarnaast wordt informatie sneller aangenomen als waarheid wanneer dit overeenkomt met persoonlijke overtuigingen \autocite{Schwalbe2024}. Deze vaststellingen benadrukken de nood aan onderzoek naar automatische detectiesystemen die journalistieke bias kunnen identificeren.

Dit systeem wordt ontwikkeld voor lezers van \gls{hln} die twijfelen aan de neutraliteit van de berichtgeving en die zich bewust willen worden van de mogelijke vooringenomenheid in hun dagelijkse nieuwsconsumptie. Het vergroten van transparantie over mogelijke bias kan bijdragen aan kritisch mediagebruik en het versterken van het vertrouwen in nieuwsmedia.

\section{\IfLanguageName{dutch}{Onderzoeksvraag}{Research question}}
\label{sec:onderzoeksvraag}

De centrale onderzoeksvraag luidt: Hoe kan een browserextensie ontwikkeld worden die artikelen van \gls{hln} automatisch classificeert op basis van hun partijdigheid, met minimale impact op de gebruikservaring?

Om deze vraag te beantwoorden, richt het onderzoek zich op de deelvragen:
\begin{itemize}
  \item Welke vormen van journalistieke bias komen voor in nieuwsartikelen?
  \item Welke linguïstische kenmerken duiden op partijdigheid?
  \item Welke bestaande methoden kunnen worden toegepast voor biasdetectie in teksten?
  \item Hoe kan de semantische betekenis van woorden behouden blijven tijdens het classificatieproces?
\end{itemize}

\section{\IfLanguageName{dutch}{Onderzoeksdoelstelling}{Research objective}}
\label{sec:onderzoeksdoelstelling}

Het doel van dit onderzoek is het ontwerpen van een \gls{poc} die automatisch journalistieke partijdigheid in nieuwsartikelen van \gls{hln} detecteert en classificeert, geïmplementeerd in een gebruiksvriendelijke browserextensie. De \gls{poc} classificeert elk artikel als \textit{neutraal} of \textit{niet-neutraal} en geeft een \textit{confidence score} weer die de zekerheid van het classificatiemodel aangeeft. Indien een artikel als \textit{niet-neutraal} wordt geclassificeerd, worden de zinnen gemarkeerd die volgens het model het meest hebben bijgedragen aan deze classificatie.

De prestaties van de \gls{poc} worden geëvalueerd aan de hand van standaard classificatiemetrics, zoals precision, recall en F1-score, berekend op een testset. Daarnaast wordt een benchmark uitgevoerd waarbij de resultaten vergeleken worden met bestaande mediabiasdetectiesystemen, zoals \textit{Ground News} en \textit{Media Bias}.

\section{\IfLanguageName{dutch}{Opzet van deze bachelorproef}{Structure of this bachelor thesis}}
\label{sec:opzet-bachelorproef}

De rest van deze bachelorproef is als volgt opgebouwd:

In Hoofdstuk \ref{ch:stand-van-zaken} wordt een overzicht gegeven van de stand van zaken binnen het onderzoeksdomein, gebaseerd op een literatuurstudie.

In Hoofdstuk \ref{ch:methodologie} wordt de methodologie toegelicht en worden de gebruikte onderzoekstechnieken besproken om de onderzoeksvragen te beantwoorden.

In Hoofdstuk \ref{ch:conclusie} wordt de conclusie gepresenteerd en een antwoord geformuleerd op de onderzoeksvragen. Daarbij wordt ook een aanzet gegeven voor toekomstig onderzoek binnen dit domein.